% Please do not change %
\documentclass[11pt]{article} % font size
\usepackage[margin=1in]{geometry} % margins
\usepackage{amsmath,amsthm,amssymb} % for the  common "math symbols"
\usepackage{algorithm}
\usepackage{algpseudocode}
\usepackage[shortlabels]{enumitem} % for enumeration
\usepackage{booktabs} % if you need tables
\usepackage{listings}
\usepackage{color}
\usepackage{quiver}
\usepackage{ulem}
\DeclareMathOperator*{\argmax}{argmax}
\DeclareMathOperator*{\argmin}{argmin}

% \theoremstyle{definition}
\newtheorem{theorem}{Theorem}[section]
\newtheorem{corollary}[theorem]{Corollary}
\newtheorem{lemma}[theorem]{Lemma}
\newtheorem{proposition}[theorem]{Proposition}
\newtheorem{conjecture}[theorem]{Conjecture}
\theoremstyle{definition}
\newtheorem{definition}[theorem]{Definition}
\theoremstyle{definition}
\newtheorem{fact}[theorem]{Fact}
\theoremstyle{definition}
\newtheorem{example/}[theorem]{Example}
\newtheorem{note}[theorem]{Note}

\newenvironment{example}
  {\renewcommand{\qedsymbol}{$\blacktriangle$}%
   \pushQED{\qed}\begin{example/}}
  {\popQED\end{example/}}

\lstset{
  frame=none,
  xleftmargin=2pt,
  stepnumber=1,
  numbers=left,
  numbersep=5pt,
  numberstyle=\ttfamily\tiny\color[gray]{0.3},
  belowcaptionskip=\bigskipamount,
  captionpos=b,
  escapeinside={*'}{'*},
  language=haskell,
  tabsize=2,
  emphstyle={\bf},
  commentstyle=\it,
  stringstyle=\mdseries\rmfamily,
  showspaces=false,
  keywordstyle=\bfseries\rmfamily,
  columns=flexible,
  basicstyle=\small\sffamily,
  showstringspaces=false,
  morecomment=[l]\%,
}
% Feel free to include additional packages (if needed), but make sure it does not override the margin/font requirements.
%
\setlength{\parindent}{0pt} % no indent for new paragraphs.
\setlength{\parskip}{5pt} % distance between paragraphs, which will be used when you have a blank line in your LaTeX code.


%%%%%%%%%%%%%%%%%%%%%%%%%%%%%%%%%%%%%%%%%%%%%%%%%%%%%%%%%%%%%%%%%%%%%%%%%%%%%%
\title{Categories, Proofs, and Processes: \\Lecture 17 Adjunctions}
%
\author{Wira Azmoon Ahmad}
\date{27 Feb 2024}
%
\begin{document}
%
\maketitle

%You may use the usual \LaTeX\ environments to structure your answers:
%\begin{align}
%        \label{basegnn}
%        \mathbf{h}_u^{(t+1)} = \sigma \bigg( \mathbf{W}_\text{self}^{(t)} \mathbf{h}_u^{(t)}  + \mathbf{W}_\text{neigh}^{(t)} \sum_{v \in N(u)} \mathbf{h}_v^{(t)} + \mathbf{b}^{(t)} \bigg): 
%\end{align}


%\begin{table}[h]
%\caption{A simple table.}
%\centering
%\begin{tabular}[t]{lrrrr}
%\toprule
%Graphs & Can & Be  & Fun & Too \\ 
%\midrule 
%$G_1$ & $1$  & $2$ & $3$ & $4$ \\ 
%$G_2$ & $-1$  & $-2$ & $-3$ & $-4$ \\ 
%$G_3$ & $0.1$  & $0.2$ & $0.3$ & $0.4$ \\ 
%\bottomrule
%\end{tabular}
%\label{tab}
%\end{table}

% A functor $F:\mathcal{C} \rightarrow \mathcal{D}$ is
% \begin{itemize}
%     \item \textbf{faithful} if each map $F_{A,B}:\mathcal{C}(A,B) \rightarrow \mathcal{D}(F\ A, F\ B)$ is injective;
%     \item \textbf{full} if each map $F_{A,B}:\mathcal{C}(A,B) \rightarrow \mathcal{D}(F\ A, F\ B)$ is surjective
% \end{itemize}


\addtocounter{section}{6}
\section{Adjunctions}
\subsection{Motivating Example}
Consider two non-isomorphic posets $(\mathbb{Z},\leq),(\mathbb{R},\leq)$ with the obvious inclusion map
\[i : (\mathbb{Z},\leq) \hookrightarrow (\mathbb{R},\leq) := z\mapsto z\] 
which preserves order 
\[z\leq z' \implies i(z)\leq i(z')\]
We can't invert $i$, but can "do our best"
Given $r\in\mathbb{R}$, find some $z$ such that $i(z)=r$.
\subsubsection{Approximate from below}
\[f(r) := \text{largest $z$ s.t. }i(z)\leq r\]
i.e. "floor of $r$". $f$ is the right adjoint of $i$. ($i$ is the left adjoint of $f$.)
\[i(z')\leq r\iff z' \leq f(r)\]
\subsubsection{Approximate from above}
\[c(r) := \text{smallest $z$ s.t. }r\leq i(z)\]
i.e. "ceiling of $r$". $c$ is the left adjoint of $i$. ($i$ is the right adjoint of $c$.)
\[r\leq i(z')\iff c(r) \leq z'\]
% https://q.uiver.app/#q=WzAsMixbMCwwLCJcXG1hdGhiYntafSJdLFsxLDAsIlxcbWF0aGJie1J9Il0sWzAsMSwiaSIsMV0sWzAsMSwiZiIsMCx7ImN1cnZlIjotMn1dLFswLDEsImMiLDIseyJjdXJ2ZSI6Mn1dLFszLDIsIiIsMCx7ImxldmVsIjoxLCJzdHlsZSI6eyJuYW1lIjoiYWRqdW5jdGlvbiJ9fV0sWzIsNCwiIiwwLHsibGV2ZWwiOjEsInN0eWxlIjp7Im5hbWUiOiJhZGp1bmN0aW9uIn19XV0=
\[\begin{tikzcd}
	{\mathbb{Z}} & {\mathbb{R}}
	\arrow[""{name=0, anchor=center, inner sep=0}, "i"{description}, from=1-1, to=1-2]
	\arrow[""{name=1, anchor=center, inner sep=0}, "f", curve={height=-12pt}, from=1-1, to=1-2]
	\arrow[""{name=2, anchor=center, inner sep=0}, "c"', curve={height=12pt}, from=1-1, to=1-2]
	\arrow["\dashv"{anchor=center, rotate=-90}, draw=none, from=1, to=0]
	\arrow["\dashv"{anchor=center, rotate=-90}, draw=none, from=0, to=2]
\end{tikzcd}\]
\begin{definition}
	In posets, the Galois connection is a special case of an adjunction.
\end{definition}
\subsection{Definition}
\begin{definition}[Adjunction, Def. 68, pg. 45]
An adjunction from $\mathcal{C}$ to $\mathcal{D}$ is a triple $(F,G,\theta)$, where functors $F,G$,
% https://q.uiver.app/#q=WzAsMixbMCwwLCJcXG1hdGhjYWx7Q30iXSxbMSwwLCJcXG1hdGhjYWx7RH0iXSxbMCwxLCJGIiwwLHsib2Zmc2V0IjotMX1dLFsxLDAsIkciLDAseyJvZmZzZXQiOi0xfV1d
\[\begin{tikzcd}
	{\mathcal{C}} & {\mathcal{D}}
	\arrow["F", shift left, from=1-1, to=1-2]
	\arrow["G", shift left, from=1-2, to=1-1]
\end{tikzcd}\]
and $\theta$ is a family of bijections
\[\theta_{A,B} : \mathcal{C}(A,G(B)) \cong \mathcal{D}(F(A), B)\]
$F$ is \textbf{left adjoint} to $G$, $G$ is \textbf{right adjoint} to $F$.
\end{definition}
\subsection{Examples}
\begin{example}
	Forgetful Functor on \textbf{Pre}, category of preoders and order-preserving maps.
	\[U: \textbf{Pre}\rightarrow \textbf{Set}\]
	\[U :: (P,\leq) \mapsto P,\quad f\mapsto f\]

\begin{note}
\textbf{Q}: Is there a canonical way to “put the order back”, such that either:
\begin{enumerate}[label=(\alph*)]
\item there will always be an arrow $(P,\leq') \rightarrow (P,\leq)$? “Approximating below”
\item there will always be an arrow $(P,\leq) \rightarrow (P,\leq')$? “Approximating above”
\end{enumerate}
\end{note}
For (a), we can let $a\leq a$ and nothing else.
\[\implies a \leq' b \implies a = b \implies 1_P(a) \leq 1_P (b)\]
$1_P$ (the function) is order preserving $\leq'$ is called the discrete order $\leq_d$.

For (b), we can let $a \leq' b$ for all $a, b \in P$
\[a \leq b \implies 1(a) \leq' 1(b)\]
$1_P$ is order preserving $\leq'$ is the indiscrete order $\leq_i$

Let $D: \textbf{Set} \rightarrow \textbf{Pre}$
\[D(A)=(A,\leq)\]
\[D(f) = f\]
\[\forall a,b, a \leq b \implies a = b \implies f(a) \leq f(a) \implies f(a) \leq f(b)\]
Any function out of a discrete order is an order preserving map $\textbf{Pre}((X, \leq ), (Y, \leq)) = \textbf{Set}(X, Y)$
\[\textbf{Pre}(D A, B) = \textbf{Set}(A, U B)\]
$D$ is the left adjoint of $U$.

Let $I: \textbf{Set} \rightarrow \textbf{Pre}$
\[I(A) = (A, \leq )\]
\[I(f) = f\]
Any function into the indiscrete preorder is order-preserving. 
\[a \leq b \implies f(a) \leq b\]
\[\textbf{Pre}((X, \leq), (Y, \leq )) = \textbf{Set}(X, Y)\]
\[\textbf{Pre}(A,IB)\cong \textbf{Set}(UA,B) \implies U\dashv I\] 
$I$ is the right adjoint of $U$.
% https://q.uiver.app/#q=WzAsMixbMCwwLCJcXHRleHRiZntQcmV9Il0sWzEsMCwiXFx0ZXh0YmZ7U2V0fSJdLFswLDEsIlUiLDFdLFsxLDAsIkQiLDIseyJjdXJ2ZSI6Mn1dLFsxLDAsIkkiLDAseyJjdXJ2ZSI6LTJ9XSxbMywyLCIiLDIseyJsZXZlbCI6MSwic3R5bGUiOnsibmFtZSI6ImFkanVuY3Rpb24ifX1dLFsyLDQsIiIsMix7ImxldmVsIjoxLCJzdHlsZSI6eyJuYW1lIjoiYWRqdW5jdGlvbiJ9fV1d
\[\begin{tikzcd}
	{\textbf{Pre}} & {\textbf{Set}}
	\arrow[""{name=0, anchor=center, inner sep=0}, "U"{description}, from=1-1, to=1-2]
	\arrow[""{name=1, anchor=center, inner sep=0}, "D"', curve={height=12pt}, from=1-2, to=1-1]
	\arrow[""{name=2, anchor=center, inner sep=0}, "I", curve={height=-12pt}, from=1-2, to=1-1]
	\arrow["\dashv"{anchor=center, rotate=-89}, draw=none, from=1, to=0]
	\arrow["\dashv"{anchor=center, rotate=-91}, draw=none, from=0, to=2]
\end{tikzcd}\]
\end{example}

\begin{example}
	% https://q.uiver.app/#q=WzAsMixbMCwwLCJcXHRleHRiZntTZXR9Il0sWzEsMCwiXFx0ZXh0YmZ7VmVjdH1fe1xcbWF0aGJie1J9fSJdLFswLDEsIlxcbWF0aGJie1J9Wy1dIiwwLHsiY3VydmUiOi0yfV0sWzEsMCwiVSIsMCx7ImN1cnZlIjotMn1dLFsyLDMsIiIsMCx7ImxldmVsIjoxLCJzdHlsZSI6eyJuYW1lIjoiYWRqdW5jdGlvbiJ9fV1d
\[\begin{tikzcd}
	{\textbf{Set}} & {\textbf{Vect}_{\mathbb{R}}}
	\arrow[""{name=0, anchor=center, inner sep=0}, "{\mathbb{R}[-]}", curve={height=-12pt}, from=1-1, to=1-2]
	\arrow[""{name=1, anchor=center, inner sep=0}, "U", curve={height=-12pt}, from=1-2, to=1-1]
	\arrow["\dashv"{anchor=center, rotate=-90}, draw=none, from=0, to=1]
\end{tikzcd}\]
	\[\mathbb{R}[-] \dashv U\]
	\[\eta_X : X\rightarrow U\mathbb{R}[X]\] 
	\[\eta_X :: e_i\in X \mapsto e_i\in U\mathbb{R}[X]\]
	$R[X]$ is the vector "spanned" by $X$.
	If $X=\{e_1,e_2,e_3\}$, then $\mathbb{R}[X] \cong \mathbb{R}^3$
	\[\mathbb{R}[f] : \mathbb{R}[X] \rightarrow\mathbb{R}[Y]\]
	\[\mathbb{R}[f] (\sum_i r_i e_i) := \sum_i r_i 	f(e_i)\]
	To find the linear map, just find where the basis elements go.
\end{example}
\subsection{Unit}
% https://q.uiver.app/#q=WzAsMixbMCwwLCJcXG1hdGhjYWx7Q30iXSxbMSwwLCJcXG1hdGhjYWx7RH0iXSxbMCwxLCJGIiwwLHsiY3VydmUiOi0yfV0sWzEsMCwiRyIsMCx7ImN1cnZlIjotMn1dLFsyLDMsIiIsMCx7ImxldmVsIjoxLCJzdHlsZSI6eyJuYW1lIjoiYWRqdW5jdGlvbiJ9fV1d
\[\begin{tikzcd}
	{\mathcal{C}} & {\mathcal{D}}
	\arrow[""{name=0, anchor=center, inner sep=0}, "F", curve={height=-12pt}, from=1-1, to=1-2]
	\arrow[""{name=1, anchor=center, inner sep=0}, "G", curve={height=-12pt}, from=1-2, to=1-1]
	\arrow["\dashv"{anchor=center, rotate=-90}, draw=none, from=0, to=1]
\end{tikzcd}\]
is an adjunction iff $\forall X \in \text{Ob}(\mathcal{C})$, there exists a \underline{unit}
\[\eta_X: X\rightarrow GF(X)\]
such that
\[\forall f: X\rightarrow G(Y), \exists ! \hat{f}: F(X) \rightarrow Y\]
where
% https://q.uiver.app/#q=WzAsMyxbMCwwLCJYIl0sWzEsMCwiR0YoWCkiXSxbMSwxLCJHKFkpIl0sWzAsMSwiXFxldGFfWCJdLFswLDIsImYiLDJdLFsxLDIsIkcoXFxoYXR7Zn0pIiwwLHsic3R5bGUiOnsiYm9keSI6eyJuYW1lIjoiZGFzaGVkIn19fV1d
\[\begin{tikzcd}
	X & {GF(X)} \\
	& {G(Y)}
	\arrow["{\eta_X}", from=1-1, to=1-2]
	\arrow["f"', from=1-1, to=2-2]
	\arrow["{G(\hat{f})}", dashed, from=1-2, to=2-2]
\end{tikzcd}\]

\begin{example}
Makes more sense as unit, unit version.
% https://q.uiver.app/#q=WzAsMixbMCwwLCJcXHRleHRiZntTZXR9Il0sWzEsMCwiXFx0ZXh0YmZ7VmVjdH1fe1xcbWF0aGJie1J9fSJdLFswLDEsIlxcbWF0aGJie1J9Wy1dIiwwLHsiY3VydmUiOi0yfV0sWzEsMCwiVSIsMCx7ImN1cnZlIjotMn1dLFsyLDMsIiIsMCx7ImxldmVsIjoxLCJzdHlsZSI6eyJuYW1lIjoiYWRqdW5jdGlvbiJ9fV1d
\[\begin{tikzcd}
	{\textbf{Set}} & {\textbf{Vect}_{\mathbb{R}}}
	\arrow[""{name=0, anchor=center, inner sep=0}, "{\mathbb{R}[-]}", curve={height=-12pt}, from=1-1, to=1-2]
	\arrow[""{name=1, anchor=center, inner sep=0}, "U", curve={height=-12pt}, from=1-2, to=1-1]
	\arrow["\dashv"{anchor=center, rotate=-90}, draw=none, from=0, to=1]
\end{tikzcd}\]
$\mathbb{R}[X]$ vector space spanned by $X$.
\[v\in \mathbb{R}[X], v=\sum_i r_i x_i,\quad r_i\in \mathbb{R}, x_i\in X\]
\[\mathbb{R}[f] : \mathbb{R} [X] \rightarrow \mathbb{R}[Y]\]
\[\mathbb{R}[f]\left(\sum_i r_i x_i\right) = \sum_i r_i f(x_i)\]
% https://q.uiver.app/#q=WzAsMTAsWzIsMSwiQiJdLFszLDEsIlUoXFxtYXRoYmJ7Un1bQl0pIl0sWzMsMiwiVShWKSJdLFswLDAsImIiXSxbNCwwLCJiIl0sWzQsNCwiZihiKSJdLFs2LDEsIlxcbWF0aGJie1J9W0JdIl0sWzYsMiwiViJdLFs3LDEsIlxcc3VtX2kgcl9pIGtfaSJdLFs3LDIsIlxcc3VtX2kgcl9pIGYoYl9pKSJdLFswLDFdLFswLDIsImYiLDJdLFsxLDIsIlUoXFxoYXR7Zn0pIiwwLHsic3R5bGUiOnsiYm9keSI6eyJuYW1lIjoiZGFzaGVkIn19fV0sWzMsNCwiIiwwLHsic3R5bGUiOnsidGFpbCI6eyJuYW1lIjoibWFwcyB0byJ9fX1dLFs0LDUsIiIsMCx7InN0eWxlIjp7InRhaWwiOnsibmFtZSI6Im1hcHMgdG8ifX19XSxbMyw1LCIiLDIseyJjdXJ2ZSI6MSwic3R5bGUiOnsidGFpbCI6eyJuYW1lIjoibWFwcyB0byJ9fX1dLFs2LDcsIlxcaGF0e2Z9IiwyLHsic3R5bGUiOnsiYm9keSI6eyJuYW1lIjoiZGFzaGVkIn19fV0sWzgsOSwiIiwyLHsic3R5bGUiOnsidGFpbCI6eyJuYW1lIjoibWFwcyB0byJ9fX1dXQ==
\[\begin{tikzcd}
	b &&&& b \\
	&& B & {U(\mathbb{R}[B])} &&& {\mathbb{R}[B]} & {\sum_i r_i k_i} \\
	&&& {U(V)} &&& V & {\sum_i r_i f(b_i)} \\
	\\
	&&&& {f(b)}
	\arrow[from=2-3, to=2-4]
	\arrow["f"', from=2-3, to=3-4]
	\arrow["{U(\hat{f})}", dashed, from=2-4, to=3-4]
	\arrow[maps to, from=1-1, to=1-5]
	\arrow[maps to, from=1-5, to=5-5]
	\arrow[curve={height=6pt}, maps to, from=1-1, to=5-5]
	\arrow["{\hat{f}}"', dashed, from=2-7, to=3-7]
	\arrow[maps to, from=2-8, to=3-8]
\end{tikzcd}\]

\end{example}

\begin{example}
Makes more sense as co-unit, unit version.

\[(-\times B) : \textbf{Set} \rightarrow \textbf{Set}\]
\[\begin{cases}
    A \mapsto A\times B\\
    f \mapsto f\times 1_B
\end{cases}
\]
\[g:X\rightarrow Y\]
\[B \Rightarrow g : B\Rightarrow X \rightarrow B \Rightarrow Y\]
\[(B\Rightarrow g)(h) := g\circ h\]
% https://q.uiver.app/#q=WzAsNixbMiwxLCJYIl0sWzMsMSwiQlxcUmlnaHRhcnJvdyhYXFx0aW1lcyBCKSJdLFszLDIsIkJcXFJpZ2h0YXJyb3cgWSJdLFswLDAsIngiXSxbNCwwLCJcXGxhbWJkYSBiLiAoeCxiKSJdLFs0LDQsImYoeCk9KFxcaGF0e2Z9XFxjaXJjIC0pKFxcbGFtYmRhIGIuKHgsYikpIl0sWzAsMSwiXFxldGFfWCJdLFsxLDIsIihmXFxjaXJjIC0pIiwwLHsic3R5bGUiOnsiYm9keSI6eyJuYW1lIjoiZGFzaGVkIn19fV0sWzAsMiwiZiIsMl0sWzMsNF0sWzMsNSwiIiwyLHsiY3VydmUiOjEsInN0eWxlIjp7InRhaWwiOnsibmFtZSI6Im1hcHMgdG8ifX19XSxbNCw1LCIiLDAseyJzdHlsZSI6eyJ0YWlsIjp7Im5hbWUiOiJtYXBzIHRvIn19fV1d
\[\begin{tikzcd}
	x &&&& {\lambda b. (x,b)} \\
	&& X & {B\Rightarrow(X\times B)} \\
	&&& {B\Rightarrow Y} \\
	\\
	&&&& {f(x)=(\hat{f}\circ -)(\lambda b.(x,b))}
	\arrow["{\eta_X}", from=2-3, to=2-4]
	\arrow["{(f\circ -)}", dashed, from=2-4, to=3-4]
	\arrow["f"', from=2-3, to=3-4]
	\arrow[from=1-1, to=1-5]
	\arrow[curve={height=24pt}, maps to, from=1-1, to=5-5]
	\arrow[maps to, from=1-5, to=5-5]
\end{tikzcd}\]
\[\hat{f} : X\times B \rightarrow Y\]
\[\hat{f}(x,b) = f(x)(b)\]
\[(\hat{f} \circ -)(\lambda b.(x,b)) = \lambda b.\hat{f}(x,b) = \lambda b.f(x)(b) = f(x)\]
Last equality by $\eta$ expansion.

\end{example}

\subsubsection{Notes}
\begin{itemize}
    \item Example: $X$ is set, $GF(X)$ is a list.
    \item More useful for the left adjoints.
\end{itemize}

\subsection{Bijection}
\[\mathcal{D}(F(X),Y) \cong \mathcal{C}(X,G(Y))\]

\subsection{Co-Unit}
% https://q.uiver.app/#q=WzAsMixbMCwwLCJcXG1hdGhjYWx7Q30iXSxbMSwwLCJcXG1hdGhjYWx7RH0iXSxbMCwxLCJGIiwwLHsiY3VydmUiOi0yfV0sWzEsMCwiRyIsMCx7ImN1cnZlIjotMn1dLFsyLDMsIiIsMCx7ImxldmVsIjoxLCJzdHlsZSI6eyJuYW1lIjoiYWRqdW5jdGlvbiJ9fV1d
\[\begin{tikzcd}
	{\mathcal{C}} & {\mathcal{D}}
	\arrow[""{name=0, anchor=center, inner sep=0}, "F", curve={height=-12pt}, from=1-1, to=1-2]
	\arrow[""{name=1, anchor=center, inner sep=0}, "G", curve={height=-12pt}, from=1-2, to=1-1]
	\arrow["\dashv"{anchor=center, rotate=-90}, draw=none, from=0, to=1]
\end{tikzcd}\]
is an adjunction iff $\forall X \in \text{Ob}(\mathcal{C})$, there exists a \underline{co-unit}
\[\varepsilon_Y: FG(Y)\rightarrow Y\]
such that
\[\forall f: F(X)\rightarrow Y, \exists ! \tilde{f}: X\rightarrow G(Y)\]
where
% https://q.uiver.app/#q=WzAsMyxbMCwwLCJGKFgpIl0sWzAsMSwiRkcoWSkiXSxbMSwxLCJZIl0sWzAsMSwiRihcXGhhdHtmfSkiLDIseyJzdHlsZSI6eyJib2R5Ijp7Im5hbWUiOiJkYXNoZWQifX19XSxbMCwyLCJmIl0sWzEsMiwiXFx2YXJlcHNpbG9uX1kiLDJdXQ==
\[\begin{tikzcd}
	{F(X)} \\
	{FG(Y)} & Y
	\arrow["{F(\tilde{f})}"', dashed, from=1-1, to=2-1]
	\arrow["f", from=1-1, to=2-2]
	\arrow["{\varepsilon_Y}"', from=2-1, to=2-2]
\end{tikzcd}\]

\begin{note}
Co-unit is some kind of evaluation.
\end{note}

\begin{example}
Makes more sense as unit, co-unit version.
\[F = \mathbb{R}[-], G= U\]
% https://q.uiver.app/#q=WzAsMixbMCwwLCJcXHRleHRiZntTZXR9Il0sWzEsMCwiXFx0ZXh0YmZ7VmVjdH1fe1xcbWF0aGJie1J9fSJdLFswLDEsIlxcbWF0aGJie1J9Wy1dIiwwLHsiY3VydmUiOi0yfV0sWzEsMCwiVSIsMCx7ImN1cnZlIjotMn1dLFsyLDMsIiIsMCx7ImxldmVsIjoxLCJzdHlsZSI6eyJuYW1lIjoiYWRqdW5jdGlvbiJ9fV1d
\[\begin{tikzcd}
	{\textbf{Set}} & {\textbf{Vect}_{\mathbb{R}}}
	\arrow[""{name=0, anchor=center, inner sep=0}, "{\mathbb{R}[-]}", curve={height=-12pt}, from=1-1, to=1-2]
	\arrow[""{name=1, anchor=center, inner sep=0}, "U", curve={height=-12pt}, from=1-2, to=1-1]
	\arrow["\dashv"{anchor=center, rotate=-90}, draw=none, from=0, to=1]
\end{tikzcd}\]
"Forget that $Y$ is a vector space, treat it as a set"
% https://q.uiver.app/#q=WzAsNSxbMCwwLCJcXG1hdGhiYntSfVtYXSJdLFsxLDEsIlkiXSxbMCwxLCJcXG1hdGhiYntSfVtVKFkpXSJdLFswLDIsIlxcc3VtX2kgcl9pIFxcY2RvdCBbdl9pXSJdLFsxLDIsIlxcc3VtX2kgcl9pIHZfaSJdLFswLDEsImYiXSxbMiwxLCJcXHZhcmVwc2lsb25fWSIsMl0sWzAsMiwiRihcXHRpbGRle2Z9KSIsMix7InN0eWxlIjp7ImJvZHkiOnsibmFtZSI6ImRhc2hlZCJ9fX1dLFszLDIsIiIsMix7InN0eWxlIjp7ImJvZHkiOnsibmFtZSI6InNxdWlnZ2x5In19fV0sWzMsNCwiIiwwLHsic3R5bGUiOnsidGFpbCI6eyJuYW1lIjoibWFwcyB0byJ9fX1dXQ==
\[\begin{tikzcd}
	{\mathbb{R}[X]} \\
	{\mathbb{R}[U(Y)]} & Y \\
	{\sum_i r_i \cdot [v_i]} & {\sum_i r_i v_i}
	\arrow["f", from=1-1, to=2-2]
	\arrow["{\varepsilon_Y}"', from=2-1, to=2-2]
	\arrow["{F(\tilde{f})}"', dashed, from=1-1, to=2-1]
	\arrow[squiggly, from=3-1, to=2-1]
	\arrow[maps to, from=3-1, to=3-2]
\end{tikzcd}\]
\[\tilde{f}(x)=f(x)\]
\end{example}

\begin{example}
Makes more sense as co-unit.
\[(-\times B) : \textbf{Set} \rightarrow \textbf{Set}\]
\[\begin{cases}
    A \mapsto A\times B\\
    f \mapsto f\times 1_B
\end{cases}
\]
\[\textbf{Set}(A\times B, C) \cong \textbf{Set}(A,B\Rightarrow C)\]
% https://q.uiver.app/#q=WzAsMixbMCwwLCJcXHRleHRiZntTZXR9Il0sWzEsMCwiXFx0ZXh0YmZ7U2V0fSJdLFswLDEsIigtXFx0aW1lcyBCKSIsMCx7ImN1cnZlIjotMn1dLFsxLDAsIkc9KEJcXFJpZ2h0YXJyb3ctKSIsMCx7ImN1cnZlIjotMn1dLFsyLDMsIiIsMCx7ImxldmVsIjoxLCJzdHlsZSI6eyJuYW1lIjoiYWRqdW5jdGlvbiJ9fV1d
\[\begin{tikzcd}
	{\textbf{Set}} & {\textbf{Set}}
	\arrow[""{name=0, anchor=center, inner sep=0}, "{(-\times B)}", curve={height=-12pt}, from=1-1, to=1-2]
	\arrow[""{name=1, anchor=center, inner sep=0}, "{G=(B\Rightarrow-)}", curve={height=-12pt}, from=1-2, to=1-1]
	\arrow["\dashv"{anchor=center, rotate=-90}, draw=none, from=0, to=1]
\end{tikzcd}\]
\[FG(C) = (B\Rightarrow C) \times B\]
% https://q.uiver.app/#q=WzAsNixbMSwzLCIoQlxcUmlnaHRhcnJvdyBDKSBcXHRpbWVzIEIiXSxbMiwzLCJDIl0sWzAsNCwiKFxcbGFtYmRhIHguZihhLHgpLGIpIl0sWzQsNCwiZihhLGIpIl0sWzEsMiwiQVxcdGltZXMgQiJdLFswLDAsIihhLGIpIl0sWzAsMSwiXFx2YXJlcHNpbG9uX0MiLDJdLFsyLDMsIiIsMCx7InN0eWxlIjp7InRhaWwiOnsibmFtZSI6Im1hcHMgdG8ifX19XSxbNCwxLCJmIl0sWzQsMCwiXFx0aWxkZXtmfSBcXHRpbWVzIDFfQiIsMix7InN0eWxlIjp7ImJvZHkiOnsibmFtZSI6ImRhc2hlZCJ9fX1dLFs1LDIsIiIsMCx7InN0eWxlIjp7InRhaWwiOnsibmFtZSI6Im1hcHMgdG8ifX19XSxbNSwzLCIiLDIseyJjdXJ2ZSI6LTEsInN0eWxlIjp7InRhaWwiOnsibmFtZSI6Im1hcHMgdG8ifX19XV0=
\[\begin{tikzcd}
	{(a,b)} \\
	\\
	& {A\times B} \\
	& {(B\Rightarrow C) \times B} & C \\
	{(\lambda x.f(a,x),b)} &&&& {f(a,b)}
	\arrow["{\varepsilon_C}"', from=4-2, to=4-3]
	\arrow[maps to, from=5-1, to=5-5]
	\arrow["f", from=3-2, to=4-3]
	\arrow["{\tilde{f} \times 1_B}"', dashed, from=3-2, to=4-2]
	\arrow[maps to, from=1-1, to=5-1]
	\arrow[curve={height=-12pt}, maps to, from=1-1, to=5-5]
\end{tikzcd}\]
\[\tilde{f}: A\rightarrow (B\Rightarrow C)\]
\[\tilde{f} :: a \mapsto \lambda x. f(a,x)\]
\[\tilde{f} :: a \mapsto f(a,-)\]
Notation mapping of the same thing.
\end{example}




\end{document}
